% webview-bridge-security.tex — WebViews as an attack surface: UIWebView vs WKWebView,
% the native <-> JS bridge (WKScriptMessageHandler, WKScriptMessage.frameInfo, reply handlers,
% WKContentWorld), injecting JS safely (callAsyncJavaScript arguments), navigation policy,
% file access, inspectability; Android addJavascriptInterface vs addWebMessageListener;
% React Native WebView (postMessage carries no verified origin, originWhitelist only
% governs navigation, onShouldStartLoadWithRequest, injectedJavaScriptObject, the 2020
% Android WebView UXSS advisory).
% Sources (checked 2026-09-25):
%   https://developer.apple.com/documentation/webkit/wkcontentworld (iOS 14; DOM shared)
%   https://developer.apple.com/documentation/webkit/wkusercontentcontroller (handler APIs)
%   https://developer.apple.com/documentation/webkit/wkwebview (callAsyncJavaScript, isInspectable)
%   https://developer.apple.com/news/?id=12232019b and ?id=edwud51q (UIWebView App Store deadlines)
%   https://github.com/react-native-webview/react-native-webview/blob/master/docs/Reference.md
%   https://github.com/advisories/GHSA-36j3-xxf7-4pqg (CVE-2020-6506, <=10.10.2, fixed 11.0.0)
%   https://developer.android.com/reference/android/webkit/WebSettings (file access defaults)
%   https://developer.android.com/reference/androidx/webkit/WebViewCompat (addWebMessageListener)
%   https://www.rfc-editor.org/rfc/rfc8252 (8.12: no embedded user-agents for OAuth)
% @labels: area=security kind=api platform=cross-platform topic=security,platform-apis discussion=247
% @tags: wkwebview, uiwebview, wkscriptmessagehandler, wkcontentworld, js-bridge, callasyncjavascript, navigation-policy, react-native-webview, postmessage, originwhitelist, addjavascriptinterface, cve-2020-6506
\documentclass[8pt]{extarticle}
\usepackage{printup-sheet}
\usepackage{array}

\lstdefinelanguage{SwiftSheet}{
  morekeywords={final,class,func,let,var,guard,else,return,async,await,try,nil,true,false,in},
  sensitive=true, morecomment=[l]{//}, morestring=[b]",
  literate={==}{{\hbox{=}\hbox{=}}}2 {!=}{{\hbox{!}\hbox{=}}}2 {->}{{\hbox{-}\hbox{>}}}2
           {===}{{\hbox{=}\hbox{=}\hbox{=}}}3 {=>}{{\hbox{=}\hbox{>}}}2 {&&}{{\hbox{\&}\hbox{\&}}}2}
\lstdefinelanguage{JSSheet}{
  morekeywords={const,if,new,true,false,return},
  sensitive=true, morecomment=[l]{//}, morecomment=[s]{/*}{*/}, morestring=[b]', morestring=[b]",
  literate={===}{{\hbox{=}\hbox{=}\hbox{=}}}3 {=>}{{\hbox{=}\hbox{>}}}2 {&&}{{\hbox{\&}\hbox{\&}}}2
           {://}{{\hbox{:}\hbox{/}\hbox{/}}}3}
\lstset{basicstyle=\ttfamily\scriptsize, aboveskip=2pt, belowskip=2pt}

\newcommand\ct[1]{\texttt{#1}}
\newcommand\css{\hbox{:}\hbox{/}\hbox{/}}   % "://" without the mono font's ligature
\newcommand\hd[1]{\par\vspace{3pt}\noindent{\bfseries\color{sheetBlue}#1}\par\vspace{1pt}}

\tikzset{
  src/.style={draw=#1, thick, rounded corners=2pt, fill=#1!7, font=\scriptsize, text width=32mm,
              align=left, inner sep=1.5pt, minimum height=6mm},
  nat/.style={draw=sheetGreen!70!black, thick, rounded corners=2pt, fill=sheetGreen!8, font=\scriptsize,
              text width=37mm, align=left, inner sep=1.5pt, minimum height=6mm},
  lbl/.style={font=\tiny, text=black!75, inner sep=1pt, align=center},
}

\begin{document}

\sheettitle{WebView \& JS-bridge security — WKWebView, React Native}{security · folio}

\oneliner{A WebView runs \emph{someone's} JavaScript inside your app. The \textbf{bridge}
(WebKit message handlers, React Native \emph{postMessage}, Android \emph{addJavascriptInterface})
hands that JavaScript native powers — and it is reachable by \textbf{every script in every frame}
of whatever page is loaded. Treat each message as untrusted input: \textbf{check the origin and
frame, allowlist commands, restrict navigation}, pass data to JS as \emph{values}, never as code.}

\vspace{2pt}
\noindent\begin{tikzpicture}[sheet]
  % web side
  \draw[sheetGrey, thick, rounded corners=4pt, fill=black!2] (0,0.45) rectangle (6.9,-3.85);
  \node[font=\bfseries\scriptsize, anchor=north west] at (0.1,0.42) {WebContent process — untrusted};
  \node[src=sheetGreen!70!black, anchor=west] (m) at (0.15,-0.25) {\textbf{main frame}: your https page};
  \node[src=sheetOrange, anchor=west] (a) at (0.15,-0.95) {\textbf{3rd-party} script / ad tag in it};
  \node[src=sheetRed, anchor=west] (x) at (0.15,-1.65) {\textbf{XSS} in your own page};
  \node[src=sheetRed, anchor=west] (f) at (0.15,-2.35) {\textbf{iframe} from another origin};
  \node[src=sheetRed, anchor=west] (n) at (0.15,-3.05) {page \textbf{navigated away} (redirect)};
  \node[lbl, anchor=north west, text=sheetRed, text width=29mm, align=left] at (3.75,0.4)
    {all five can call \ct{messageHandlers.app}\\\ct{.postMessage(…)} unless you\\check the sender};
  % the gate
  \node[draw=sheetBlue, very thick, fill=sheetBlue!10, rounded corners=3pt, text width=34mm, align=left,
        font=\scriptsize, inner sep=2pt] (g) at (8.95,-1.45)
    {\textbf{bridge checkpoint}\\
     1 \ct{frameInfo.isMainFrame}\\
     2 \ct{securityOrigin}: https + your host\\
     3 decode into a typed \ct{Command}\\
     4 allowlist: \emph{this} page may do \emph{this}\\
     5 sensitive? user confirms / re-auth\\
     \textcolor{sheetRed}{RN \ct{onMessage}: a bare string —}\\
     \textcolor{sheetRed}{no frame/origin to check}};
  \foreach \s in {m,a,x,f,n} \draw[hot, draw=sheetGrey] (\s.east) -- (g.west);
  % native side
  \node[font=\bfseries\scriptsize, text=sheetGreen!50!black, anchor=west] at (11.2,0.3) {native — your privileges};
  \node[nat, anchor=west] (k) at (11.2,-0.25) {Keychain tokens, session};
  \node[nat, anchor=west] (p) at (11.2,-0.95) {payments, contacts, camera};
  \node[nat, anchor=west] (fs) at (11.2,-1.65) {files, \ct{openURL} to other apps};
  \node[nat, anchor=west] (c) at (11.2,-2.35) {navigation: \ct{decidePolicyFor}};
  \node[nat, anchor=west, draw=sheetBlue, fill=sheetBlue!6] (j) at (11.2,-3.05) {native $\to$ JS: \ct{callAsyncJavaScript} with \ct{arguments}};
  \foreach \t in {k,p,fs} \draw[flow, draw=sheetGreen!60!black] (g.east) -- (\t.west);
  \draw[flow, draw=sheetBlue] (j.west) -- ++(-0.35,0) |- node[lbl, pos=0.75, below]{values, never string concatenation} (6.9,-3.45);
  \node[lbl, anchor=north west, text=sheetGrey, align=left] at (0.1,-3.95)
    {\textbf{Isolation}: a \ct{WKContentWorld.world(name:)} gives your injected script its own JS globals + handlers the page can't see —
     but the \textbf{DOM is shared}, so a named world hides \emph{your} code, it does not make the page trustworthy.};
\end{tikzpicture}

\begin{multicols}{2}
\raggedright\setstretch{1.0}

\hd{iOS}
\begin{itemize}
  \item \textbf{UIWebView} (deprecated iOS 12) ran web content \emph{in your process}; the App
        Store refused new apps using it from Apr 2020. \textbf{WKWebView}: separate WebContent
        process, Safari's engine + JIT.
  \item Bridge: \ct{userContentController.add(handler, name: "app")} $\to$ JS
        \ct{window.webkit.messageHandlers.app.postMessage(obj)} (JSON-like types). The handler gets
        a \ct{WKScriptMessage}: \ct{body}, \ct{frameInfo} (\ct{isMainFrame}, \ct{request},
        \ct{securityOrigin}), \ct{world}. Reply: \ct{WKScriptMessageHandlerWithReply} (iOS 14) —
        JS gets a Promise.
  \item \textbf{Content worlds} (iOS 14): \ct{.page}, \ct{.defaultClient}, \ct{.world(name:)};
        \ct{add(\_:contentWorld:name:)} exposes a handler \emph{only} in that world.
  \item \textbf{Native $\to$ JS}: \ct{evaluateJavaScript("show('\textbackslash(name)')")} is
        injection (a quote in \ct{name} ends the string).
        \ct{callAsyncJavaScript(\_:arguments:in:contentWorld:)} passes \emph{values}.
  \item \textbf{Navigation}: \ct{decidePolicyFor} — allow https + your hosts; others $\to$
        Safari; cancel \ct{javascript:}, \ct{file:}, \ct{data:}, custom schemes. iOS 14
        \ct{WKAppBoundDomains} + \ct{limitsNavigationsToAppBoundDomains} disable script
        injection / handlers elsewhere.
  \item \textbf{Files}: \ct{loadFileURL(\_:allowingReadAccessTo:)} — the narrowest directory,
        never the whole container. \textbf{Debug}: \ct{isInspectable} (iOS 16.4, default false) —
        keep it off in release.
  \item The controller \textbf{retains} its handlers $\to$ handler $\to$ VC $\to$ WebView cycle:
        weak proxy, or \ct{removeScriptMessageHandler(forName:)}.
\end{itemize}

\hd{Android}
\ct{addJavascriptInterface}: every \ct{@JavascriptInterface} method, to every frame and origin
(before API 17: reflection $\to$ RCE). Prefer AndroidX \ct{WebViewCompat.addWebMessageListener}
with \ct{allowedOriginRules} — the listener gets \ct{sourceOrigin} + \ct{isMainFrame}.
\ct{setAllowFileAccess} defaults \ct{false} from API 30. \ct{shouldOverrideUrlLoading}
= \ct{decidePolicyFor}.

\hd{Example — a checked bridge (Swift) and RN}
\begin{lstlisting}[language=SwiftSheet]
final class Bridge: NSObject, WKScriptMessageHandlerWithReply {
  func userContentController(_ c: WKUserContentController,
      didReceive m: WKScriptMessage) async -> (Any?, String?) {
    let o = m.frameInfo.securityOrigin
    guard m.frameInfo.isMainFrame, o.protocol == "https",
          o.host == "app.example.com",
          let cmd = Command(json: m.body)    // typed, allowlisted
    else { return (nil, "denied") }
    return (await perform(cmd), nil)
  }
}
cfg.userContentController.addScriptMessageHandler(Bridge(),
    contentWorld: .page, name: "app")
try await web.callAsyncJavaScript("render(item)",
    arguments: ["item": title], contentWorld: .page)
\end{lstlisting}
\begin{lstlisting}[language=JSSheet]
<WebView source={{ uri: 'https://app.example.com' }}
  originWhitelist={['https://app.example.com']} // navigation only
  onShouldStartLoadWithRequest={r =>
    new URL(r.url).origin === 'https://app.example.com'}
  onMessage={e => { const m = parseCommand(e.nativeEvent.data)
                    if (m && allowed(m)) run(m) }}   // string in
  javaScriptEnabled setSupportMultipleWindows /* keep true */ />
\end{lstlisting}

\hd{React Native WebView — the traps}
\begin{itemize}
  \item \ct{onMessage} receives \ct{window.ReactNativeWebView.postMessage(string)} — a
        \textbf{string only}, with no verified sender: any script the page runs can post.
  \item \ct{originWhitelist} (default \ct{http\css{}*}, \ct{https\css{}*}) lists origins the WebView
        may \textbf{navigate} to; others are handed to the OS. It does \emph{not} filter messages.
  \item \ct{onShouldStartLoadWithRequest}: \ct{false} blocks; not called on Android's first load.
  \item \ct{injectedJavaScriptObject} is visible to \emph{all} frames;
        \ct{injectedJavaScriptForMainFrameOnly} defaults \ct{true}.
  \item \textbf{GHSA-36j3-xxf7-4pqg / CVE-2020-6506} (Sep 2020, moderate): Android WebView before
        83.0.4103.106 let a cross-origin \textbf{iframe} run code in the top document (UXSS);
        \ct{react-native-webview} $\leq$ 10.10.2 exposed, 11.0.0 mitigates —
        \ct{setSupportMultipleWindows=false} re-opens it.
\end{itemize}

\hd{Traps}
\begin{itemize}
  \trap{``It's our own page'' — until a CDN script, an ad or one XSS runs on it.}
  \trap{Login / OAuth in an embedded WebView: the app can read the password;
        RFC 8252 says use the system browser (\ct{ASWebAuthenticationSession}).}
  \trap{Injecting the refresh token into JS for the web part — any page script can read it.}
\end{itemize}

\hd{Remember}
\textbf{``Who sent it, may they, is it well-formed?''} — then values out, never code.


\end{multicols}

\noindent{\footnotesize\color{sheetGrey}\textit{Related:} owasp-masvs-mobile-top10 (PLATFORM-2, M4) ·
deep-link-attack-surface · oauth-oidc-jwt · rn-native-modules · mobile-data-leakage (WebView store)}

\end{document}
